\newpage

\chapter*{Everyone Is Their Own CEO}

\markboth{Everyone Is Their Own CEO}{}
\phantomsection
\addcontentsline{toc}{chapter}{Everyone Is Their Own CEO}
\label{ch:llm_ceo}

% Custom commands for formatting
\newcommand{\interviewQ}[1]{%
    \vspace{0.8em}
    \noindent\textbf{Q: #1}
    \par
}

\newcommand{\interviewA}[1]{%
    \vspace{0.4em}
    \noindent\textbf{Xun Jiang: }#1
    \par
}

% Editor's note
\begin{thinkbox}[Editor's Note:]
With the widespread adoption of generative AI and related technologies, a wave of "Super Individuals" has emerged---people who can wear multiple hats by leveraging AI tools. These employees have dramatically expanded the boundaries of their personal capabilities, often completing tasks independently that previously required coordination across multiple roles. This trend is driving a profound transformation in organizational management. At the Q2 Group Strategy Roundtable in 2025, our chairman and the heads of various innovation ventures shared ideas and proposed building "an ecological rainforest for Super Individuals."

In this edition, the editorial team interviewed Xun Jiang, CEO of Theta, to learn about his understanding and practice of the Super Individual concept. Since April 2025, Theta has been actively implementing management reforms to maximize the value of Super Individuals, and several promising early results have already emerged.
\end{thinkbox}

% Main text begins

\section{What Does a Super Individual Look Like?}
\interviewQ{What is the difference between a Super Individual and what we traditionally call a "generalist"?}

\interviewA{The two are fundamentally different. A traditional "generalist" or "full-stack engineer" typically refers to someone with comprehensive technical skills who can handle both front-end and back-end work and solve a wide range of problems. But they tend to wait passively for task assignments and rely on their technical abilities to address whatever comes their way.

A "Super Individual," on the other hand, operates from a fundamentally different logic---starting from business problems. Super Individuals are not just technically versatile; they possess integrated capabilities spanning business, product, operations, and more. They can identify concrete business opportunities and independently handle everything from understanding user needs and designing solutions to writing code, deploying, and driving adoption.

\newpage
\textbf{In} the past, building a product required writing detailed PRDs, going through round after round of discussions, and navigating lengthy, complex processes. Now, with AI tools like GPT, a Super Individual can rapidly turn ideas into prototypes and validate whether a business opportunity holds---sometimes producing an initial product in just one or two days. Once the prototype exists, communication costs plummet. Everyone can see it immediately, and a small group can quickly rally to push things forward.

\vspace{1em}

The empowerment that Super Individuals gain goes far beyond the technical level. We are now in a problem-driven paradigm: as long as someone identifies a concrete business problem and has the self-motivation to solve it, AI can massively amplify their personal capabilities, enabling them to solve real problems end to end. Such people have sharper insight into societal and user needs, no longer constrained by limited resources or team size. Many ideas that were once impossible to execute due to individual limitations can now become reality---all it takes is the courage to imagine and the willingness to try.

\vspace{1em}

In short, the biggest difference with Super Individuals lies in their intense business drive and rapid validation capability---the ability to efficiently seize opportunities and execute on them, rather than merely being a well-rounded "generalist."}

\interviewQ{At the group strategy meeting, it was mentioned that teams led by Super Individuals have characteristics of being "high-caliber, small, and lean." How do you see these three qualities? Or what do you think is the most important trait of a Super Individual?}

\interviewA{In my view, the most essential trait of a Super Individual is the business acumen to identify real user pain points.

First, they can keenly detect a specific, nuanced user need and continuously ask themselves: as an ordinary user, does this problem also trouble many others?

Second, they have strong self-motivation and are willing to personally tackle the problem. They mentally assess whether they can handle it alone, or whether they need to assemble a lean team of two or three people with complementary skills. Typically, a five-person team can already cover most key roles---back-end, front-end, marketing, and so on. Each person brings their own specialty, and with the added power of AI tools like GPT, the entire team's capability becomes formidable.

\vspace{1em}

Therefore, I see the core of a Super Individual as the ability to define a concrete problem and have the confidence to solve it personally. They can build prototypes extremely quickly, and once a prototype exists, communication costs drop to near zero. People stop debating in the abstract and start looking at actual systems and products, which attracts more people to join, creating a snowball effect. The Windsurf team, for example, reached tens of millions in ARR within eight months---a perfect illustration of this logic.

\textbf{Regarding} opportunities for startups in the AI era, there are two main schools of thought. One holds that small companies have limited opportunity because AI is fundamentally an efficiency-enhancing tool; large companies already have established user bases and can quickly dominate the market by embedding AI.

The other view observes that many small teams today (like OpenEvidence, Cursor, etc.) have only about twenty people, yet they identify problems and solve them in entirely new ways. These companies deserve our serious study. Their operating styles may be unconventional and their business models debatable, but they have already generated significant traction. Some may get acquired or disappear, but they have survived for two years and changed how people understand and use search.

\vspace{1em}

Recently, the search industry is being restructured, especially in standardized Q\&A and information retrieval scenarios. An increasing number of users are turning to large language models first, which shows that change is happening much faster than we imagined. The first things to be disrupted are often not entire professions themselves, but rather the standardized information gathering, initial screening, and clerical components within them.

\vspace{1em}

The key lies in one's understanding of how the world will change over the next few years. Once you possess this understanding, today's technology is entirely capable of supporting a very small team to rapidly turn insights into prototypes, and even generate revenue relatively quickly, despite limited resources.

\vspace{1em}

Following this logic, corporate structures will evolve accordingly. In the past, the emphasis was on teamwork---you needed large teams and ample resources to get things done. But now, with Super Individuals, the logic has changed. You have a "machine colleague" that understands you and communicates with you seamlessly, dramatically boosting efficiency. The communication cost between humans is extremely high, and complexity grows quadratically with headcount. In the past, success required stacking people and capital to build teams. But now, with tools like GPT, individuals become highly autonomous---they only need to collaborate with AI, and the friction is minimal.

\vspace{1em}

You can mentally rehearse and simulate the entire business loop in your head---from users, product, to operations---constantly thinking through problems. When you encounter something you don't know, you consult GPT, which fills in your knowledge and capabilities in real time.
A single person can quickly run through the entire loop, and the efficiency gain is staggering. Everyone operates like a CEO, capable of independently completing the journey from idea to prototype. This is especially valuable in the startup phase. Once the business takes shape and the team grows, coordination still introduces some efficiency loss, but the speed from idea to prototype is unprecedented.

% Theta's current approach to identifying and cultivating Super Individuals is also based on the "big data business" direction confirmed by our chairman. Within this framework, we encourage everyone to independently define core business problems and solve them end to end. We no longer divide work by job title, but by functional module---one person or a small team is responsible for a complete business module. From front-end to back-end, from page development to deployment, everything can be done independently, with questions directed to colleagues or GPT as they arise. Integration used to be a headache; now one person can do the "integration" in their head, and friction is drastically reduced.

\textbf{In} summary, a Super Individual is someone who can rapidly close the loop around a specific problem, possesses strong self-motivation, and excels at learning and communication---whether with people or with AI. Future company structures will be leaner and more agile, requiring less capital, with businesses that can sustain themselves. As long as you have ideas and the will to act, resources will rarely be the limiting factor. What people are really competing on is their insight into the future. The business of the future is about earning "consensus money"---whoever can discover and realize consensus ahead of others captures the value.}

\interviewQ{The company has been operating with five-person small teams for a while now. Can you describe how this works in practice?}

\interviewA{We officially rolled out the small team model around late April or early May this year, with the key catalyst being the release of Claude 3.7, which brought new breakthroughs in AI capability. At the time, our business direction wasn't fully defined, but the strategic direction was set. A team proposed building a to-provider product, so we assembled a standard five-person team: one product manager, one front-end developer, one UI designer, and two AI engineers. With functional coverage in place, team members didn't confine themselves to their defined roles---instead, they strived to become "Super Individuals," collaborating to drive the entire process from requirement definition to end-to-end execution.

% This team incubated a product in two to three months and is about to sign the first contract in North America. From the very beginning, we established a rule: team size cannot expand arbitrarily. We only consider adding members when revenue can cover the cost of new hires. The company can provide initial support, but afterward the team must sustain itself. Unless there are special circumstances like a large-scale marketing push, the basic principle is self-sufficiency.

\vspace{1em}

Another example is that I personally rewrote the MCP service from scratch within two to three weeks---setting aside daily management duties, closing the loop entirely on my own, thinking independently, and doing the hands-on work from start to finish. With zero communication overhead, efficiency was extremely high; problems were solved immediately, and my thinking remained crystal clear. I documented this thought process in the article "Along the Value Stream: What Will Apps Look Like When LLM-OS Arrives."

% This was actually a pivotal moment in the company's business transformation. We now position ourselves as a pure data company---health is merely our entry point into a high-density data scenario. At its core, we focus on "dirty data in, clean data out," doing nothing but data processing, and no longer chasing the competition around models and agents.

% This transformation combined our own experience with industry trends, projecting the future shape of the data business. Everyone is fighting over the agent track, but data processing is the core capability, and we aim to become the data foundation for all models and agents. The goal is crystal clear: when models like GPT open up to hundreds of millions of users, we can capture the traffic dividend with the best brand and product. Users just upload their raw health data, and after our processing, the quality of every model's responses improves.

% All current efforts revolve around polishing our product according to the first principle of "dirty data in, clean data out."

With first principles as our guide, everyone can become a Super Individual with a clear sense of direction. Take Celine, for example. As a product manager, she can now independently develop applications on her phone---her development capability even surpasses mine. We hope that through the power of role models, everyone will come to believe: things that were impossible before can now be done by a single person. All it takes is the courage to imagine and try---practice will prove that individual potential far exceeds what anyone imagined.

% Everyone can focus on a small entry point, dig deep, and pursue excellence. The requirement is simple: stop saying you need integration meetings. Our codebase is now a unified monorepo where everyone can view and modify all the code. Front-end and back-end are fully interconnected; anyone can modify any part. The tech stack has also been unified---Node.js for the front-end, Python for the back-end---lowering the learning and maintenance costs. Problems can be solved independently without depending on others.

% Therefore, Theta has always kept the team size small and has no plans to expand in the near term. We firmly believe that small teams and high efficiency are the optimal model for future innovation.

}

\interviewQ{Are people with "Super Individual" qualities born that way, or can they be developed?}

\interviewA{Some people are indeed natural-born Super Individuals, but quite a few have been cultivated over time. For instance, one engineer on our team originally did only back-end development and was initially struggling with front-end code. Recently, he came to me and said: "Hey Jiang, I wrote a front-end page---take a look." I checked it out and thought it was excellent.

\textbf{If} all the code can be deployed locally, then anyone can run and read the entire codebase, which eliminates all collaboration friction. We are still at a very early stage---the business hasn't really taken off yet. There's no need to obsess over division of labor and processes; we can worry about those things once the business succeeds. Besides, at our iteration speed, we can rewrite the entire codebase in three months---sometimes even in two weeks---so many traditional constraints can simply be broken through.

\vspace{1em}

Another approach we're pursuing is bringing in more external developers. The core of the Super Individual model isn't about shrinking teams down to one or two people---it's about a shift in mindset: more open, more business-driven. The key is approaching things with a problem-solving mentality and completing the work end to end. Whether it's an internal employee or an external developer, once a problem is defined and validated within our framework, they can independently solve the entire problem. We are also considering open-sourcing some of our data and APIs---first building out the foundational information and interfaces, then providing plugin-like extension points. External developers can follow our guidelines and independently implement everything from requirement definition to data retrieval, end to end.

The key to the Super Individual lies in solving problems end to end, not in reducing headcount. It's this capability and mindset that we are truly pursuing.
}

\section{How to Become a Super Individual}

\interviewQ{What challenges have you encountered while implementing the Super Individual model?}

\interviewA{The core challenge has been keeping the team focused. Especially around April and May, as a startup, splitting the team into very small units created real pressure. A startup scattering its people so thinly didn't seem to make logical sense.

Looking back, we've moved past that initial exploration phase and entered a new stage of development. Everyone went out to try things, explore, and think, and eventually we regrouped---the entire team aligned on a shared vision and focused on one thing.
Each person on the team now owns a specific module---devices, files, agents, CDA, etc.---and is responsible for their part end to end. These modules all belong to the same overarching framework. You need to propose your own solutions, tell everyone who you need to collaborate with, and then everyone divides up the work. That's the overall logic of how we move forward.}

\interviewQ{Have there been cases where a Super Individual's strong autonomy led them in a direction misaligned with the overall strategy?}

\interviewA{In the early days, this did happen. Everyone is pursuing ultimate success and weighing the probability of that success, so in the initial phase, people might think, "I'm more reliable than the others---I should go it alone." But this phase is actually important: everyone rapidly thinks through the problem end to end, mapping out all the issues, and eventually converges on the same goal. For example, when we decided to focus on data---turning dirty data into clean data---everyone bought in, because after deep reflection, they recognized that this path had the highest probability of success.

\vspace{1em}

If someone genuinely disagrees, they should be encouraged to try their own approach. Let them lead a small team; if resources are insufficient, they can request support from management. As long as the opportunity is evaluated as credible, it should be allowed to operate independently.

% In fact, we have this kind of division right now. Currently there are two clear business directions: the core team focused on data, and a 5-person team building a to-doctor product on top of the data business. There is no third project. Newly incubated lean teams must align with the overall direction---they need to both leverage and contribute back to the core. For example, the to-doctor team builds its product on our data services; otherwise, a 5-person team competing against large companies in the market would have no chance of survival and would just be wasting time. Now, building on our data capabilities, they are creating a differentiated doctor service, which I believe has real potential. If the doctor service gains traction, it could also feed back into the data business through data collection and user growth.

}

% \interviewQ{Based on Theta's experience, do you think this model can be applied to other business teams?}

% \interviewA{The two-person or one-person company model is already very popular in the US, and I believe it is completely viable. Ultimately, teams should not grow too large---the ideal size converges to between 5 and 20 people. Being too large or too dispersed is detrimental to efficient collaboration. I am very confident in this model of small teams highly focused on one thing. If a team has ten new ideas every day and constantly changes direction, that may be the mode of a research institute, but it's not suitable for a company. For a company, the endgame is always focusing on one clear objective.

% A company cannot spin off two or three people to start a new project every time a new idea emerges---that's a recipe for failure. Therefore, I believe that only when the team is focused on the core business can selected members be allowed to build differentiated products they "believe in" on top of that core---this may be the viable path from "believing" to "seeing."

% Theta currently has four open headcount positions. If in the future someone discovers a new opportunity that synergizes with our core business and wants to pursue it, I wouldn't completely rule out that possibility.

% But given Theta's current state, expanding the team further would only introduce more friction costs. Until we have stronger profitability, I want to keep the team lean.}

\interviewQ{In the process of implementing the Super Individual and small team practices, what has been your greatest personal takeaway?}

\interviewA{The greatest takeaway has been making everyone more willing to proactively think about the future. I have always believed that success stems from breakthroughs in understanding. Whoever can discover future consensus first can seize the opportunity and act on it. It was precisely this kind of understanding that led us to discover an opportunity that hasn't yet attracted widespread attention, and we committed firmly to this path. We won't change direction easily until we either hit a dead end or achieve success.

\vspace{1em}

Whether it's the original team leaders or other members, someone on the team must step up---true "locomotives" need to emerge. These locomotives must vigorously exchange ideas, ultimately forging internal consensus within the team. Everyone rallies behind a particular locomotive's direction and focuses on the same thing.

\vspace{1em}

The Matthew effect will become increasingly pronounced in the future, and teams must be intensely focused. When someone sees an opportunity, they must go all in to make it succeed. In the early stages, we help these "seeds" sprout and let them compete fully against one another. You can even have several people pursue the same direction simultaneously, then ultimately back the strongest one and concentrate all resources on them. Either succeed or fail---you must approach it with an all-or-nothing mentality.}
